\documentclass{template}

\usepackage{mathtools}
\pgfplotsset{compat=1.18}

\newcommand{\qedwhite}{\hfill \ensuremath{\Box}}


\begin{document}

\begin{center}
    \Large \textcolor{blue!70!black}{\textbf{Commutative Algebra Problem Set \#2}} \\
    \normalsize \textit{Dmytro Zakharov}
\end{center}

\begin{problem}[Exercise 1]
    Consider the following rings:
    \begin{align*}
        &A_1 = \mathbb{C}[x] && &&A_2 = \mathbb{C}[x,x^{-1}] && &&A_3 = \mathbb{C}[x^2,x^3] \\
        &A_4 = \mathbb{C}[x,y] && &&A_5 = \mathbb{C}[x]/\langle x^6 \rangle && &&A_6 = \mathbb{C}[x,y]/\langle x^2, y^3 \rangle \\
        &A_7 = \mathbb{C}[x,y]/\langle xy \rangle && &&A_8 = \mathbb{C}[x,y] / \langle xy^2 \rangle && && A_9 = \mathbb{C}[x,y] / \langle xy, y^2 \rangle
    \end{align*}
    Prove that the rings $A_i$ and $A_j$ are not isomorphic as
$\mathbb{C}$-algebras for $i \neq j$. Can you show they are also not isomorphic
as rings?
\end{problem}

\textbf{Solution.} 

\textcolor{blue!70!black}{\textbf{\underline{$A_2$.}}} Now let us look at $A_2$.
It is clearly not isomorphic to the rest due to the nature of invertible
elements $A_2^{\times}$. Indeed, $A_2^{\times}$ consits of all elements of form
$cx^n$ for $c \in \mathbb{C}$ and $n \in \mathbb{Z}$. All other rings have units
in $\mathbb{C}^{\times}$ only.

\textcolor{blue!70!black}{\textbf{\underline{$A_3$.}}} This ring consists of all
finite sums $\sum_i\sum_j c_{i,j}x^{2i+3j}$ for $i,j \in \mathbb{Z}_{\geq 0}$.
Note that every integer $n \neq 1$ can be written as $n=2i+3j$ for some $i,j \in
\mathbb{Z}_{\geq 0}$\footnote{Indeed, denote by $S=\{2i+3j: i,j \in
\mathbb{Z}_{\geq 0}\}$. Clearly, $2 \in S$. Also note that if $y \in S$, then
$y+1 \in S$: if $y=2i_0+3j_0$, then $y+1=2(i_0-1)+3(j_0+1)$. Thus,
$S=\mathbb{Z}_{\geq 0} \setminus \{1\}$.}. So we conclude that $A_3$ consists of
polynomials $\sum_j a_jx^j$ with $a_1=0$. What are some properties of $A_3$ that
are not present, say, in $\mathbb{C}[x]$? First, it is not a UFD: for instance, $x^6$ has
two distinct decompositions: $(x^2)^3$ and $(x^3)^2$. In particular, it is not a
PID (since UFD implies PID). Moreover, $A_3 \cong \mathbb{C}[x,y]/\langle
y^2-x^3 \rangle$ (where isomorphism is given by $\varphi: x \mapsto x^2, y
\mapsto x^3$)\footnote{Credits to
\href{https://math.stackexchange.com/a/5081023}{this \texttt{mathstackexchange}}
thread.}, and the latter gives a slighly more interesting intuition of $A_3$.
For instance, the curve $C=\{(x,y) \in \mathbb{C}^2: y^2=x^3\}$ is singular at
$(0,0)$. In particular, maximal ideals are in 1-to-1 correspondence with points
on $C$: if, say, $(\alpha,\beta) \in C$, then $\langle x-\alpha, y-\beta
\rangle$ is a maximal ideal. Thus we observe that every maximal ideal of $A_3$
is generated by two elements at least, compared to other rings $A_i$.

\textcolor{blue!70!black}{\textbf{\underline{$A_5$ and $A_6$.}}} The ring $A_5$
is isomorphic to the ring of truncated polynomials $\mathbb{C}[x]_{<6}$. The
dimension of $A_5$ (as $\mathbb{C}$-vector space with basis $\{1,x,\dots,x^5\}$)
is $6$, so it is clearly not isomorphic to anything else besides probably $A_6$
(which also has a finite dimension of $6$ with the basis given by
$\{1,x,y,y^2,xy,xy^2\}$). However, it is not isomorphic to $A_5$ since maximal
ideal $\mathfrak{m}$ of $A_5$ is generated by a single element $x$, while in
$A_6$ the ideal $\mathfrak{m}/\mathfrak{m}^2$ (if we vanish all quadratic terms
$x^2,xy,y^2$) is a vector space over two elements $x$ and $y$.

\textcolor{blue!70!black}{\textbf{\underline{$A_7$, $A_8$, and $A_9$.}}} We
consider nilradicals of these rings. Note that $\text{nil}(A_7)=0$,
$\text{nil}(A_8) = \langle xy \rangle$ and $\text{nil}(A_9) = \langle y
\rangle$. We see immediately that $A_7$ and $A_8$ are not isomorphic because
$A_7$ is reduced while $A_8$ is not. Now compare quotients
$A_8/\text{nil}(A_8)$, and $A_9/\text{nil}(A_9)$. Notice that
$A_8/\text{nil}(A_8) = \mathbb{C}[x,y] / \langle xy^2 \rangle / \langle xy
\rangle \cong \mathbb{C}[x,y]/\langle xy \rangle$. In turn, $A_9/\text{nil}(A_9)
= \mathbb{C}[x,y]/ \langle xy, y^2 \rangle / \langle y \rangle \cong
\mathbb{C}[x,y]/\langle xy,y^2,y \rangle \cong \mathbb{C}[x,y]/\langle xy,y
\rangle \cong \mathbb{C}[x]$. Since $\mathbb{C}[x,y]/\langle xy \rangle$ is not
an integral domain (take $x,y \in \mathbb{C}[x,y]/\langle xy \rangle$ and note
that their product is $0$), $\mathbb{C}[x] \not\cong \mathbb{C}[x,y]/\langle xy
\rangle$ and thus $A_8$ is not isomorphic to $A_9$.

\begin{problem}[Exercise 2]
    Let $X = \text{Spec}(A)$ and $x,y \in X$. Show that:
    \begin{enumerate}
        \item[\textit{(a)}] $x$ is a closed point (i.e. $\{x\} \subseteq X$ is a
        closed subset) iff the corresponding prime ideal $\mathfrak{p}_x
        \subseteq A$ is maximal;
        \item[\textit{(b)}] the closure $\overline{\{x\}} = V(\mathfrak{p}_x)$;
        \item[\textit{(c)}] $y \in \overline{\{x\}}$ iff $\mathfrak{p}_x \subseteq \mathfrak{p}_y$.
    \end{enumerate}
\end{problem}

\textbf{Solution.} 

\underline{\textit{Part (a)}}. Suppose $\mathfrak{p}_x$ is maximal. Then by definition there
is no proper prime ideal containing $\mathfrak{p}_x$. Therefore,
$V(\mathfrak{p}_x) = \{\mathfrak{p}_x\} = \{x\}$, and so $\{x\}$ is in fact
closed. Now suppose $\{x\}$ is closed, then $\{x\} = V(\mathfrak{p}_x)$. Suppose
$\mathfrak{p}_x$ is not maximal, meaning there exists some $\mathfrak{p}_y$
between $\mathfrak{p}_x$ and $A$. But then $\mathfrak{p}_y \in
V(\mathfrak{p}_x)$ and so $V(\mathfrak{p}_x)$ contains at least two elements.
This is a contradiction. \qedwhite

\underline{\textit{Part (b)}}. We prove the auxilary proposition.

\begin{propositions}[Atiyah-Macdonald 1.11, (ii)]
    Let $\mathfrak{a}_1,\dots,\mathfrak{a}_n$ be ideals and $\mathfrak{p}
    \supseteq \bigcap_{i=1}^n \mathfrak{a}_i$. Then, $\mathfrak{p} \supseteq
    \mathfrak{a}_i$ for some $i$. In particular, if $\mathfrak{p} =
    \bigcap_{i=1}^n \mathfrak{a}_i$, then $\mathfrak{p}=\mathfrak{a}_i$.
\end{propositions}

\textbf{Proof.} We prove the contrapositive statement. Suppose $\mathfrak{p}
\not\subseteq \mathfrak{a}_i$ for all $i$. Then, there exist $x_i
\in \mathfrak{a}_i \setminus \mathfrak{p}$. Now note $\prod_{i=1}^n x_i \in
\prod_{i=1}^n \mathfrak{a}_i \subseteq \bigcap_{i=1}^n \mathfrak{a}_i$, but
$\prod_{i=1}^n x_i \not\in \mathfrak{p}$ since $\mathfrak{p}$ is prime. Hence
$\mathfrak{p} \not\supseteq \bigcap_{i=1}^n \mathfrak{a}_i$. \qedwhite

Now, \textit{part (b)} is an immediate consequence. By definition,
$\overline{\{x\}} = \bigcap_{x \in V(I)}V(I)$. Since closure is a closed set,
$\overline{\{x\}}$ is a prime ideal, so $\{x\} = V(I)$ for some ideal $I$. This
is exactly what we denote by $\mathfrak{p}_x$. \qedwhite

\underline{\textit{Part (c)}}. If $y \in \overline{\{x\}}$, then $\{y\} \subseteq
V(\mathfrak{p}_x)$. Take closure of both sides. One then has $V(\mathfrak{p}_y)
\subseteq V(\mathfrak{p}_x)$, implying $\mathfrak{p}_x \subseteq
\mathfrak{p}_y$. We can do the same in the other direction. If $\mathfrak{p}_x
\subseteq \mathfrak{p}_y$, then $\overline{\{y\}} = V(\mathfrak{p}_y) \subseteq
V(\mathfrak{p}_x) = \overline{\{x\}}$, which implies $y \in \overline{\{x\}}$.
\qedwhite

\begin{problem}[Exercise 3]
    We say that topogical space $X$ is \textbf{irreducible} if for every pair of
    closed subsets $Y_0,Y_1 \subseteq X$ such that $X = Y_0 \cup Y_1$, we have
    $X=Y_0$ or $X=Y_1$. Show that $\text{Spec}(A)$ is irreducible if and only 
    if the nilradical $\sqrt{0} \subseteq A$ is a prime ideal.
\end{problem}

\textbf{Proof.} Recall that set $Y$ is closed if $Y=V(I)$. Now, if
$\text{Spec}(A) = Y$, then $V(I) = \text{Spec}(A)$ or equivalently $I \subseteq
\bigcap_{\mathfrak{p} \in \text{Spec}(A)}\mathfrak{p}$. According to the
previous lectures, $\bigcap_{\mathfrak{p} \in \text{Spec}(A)}\mathfrak{p} =
\sqrt{0}$. Therefore, $V(I)=\text{Spec}(A)$ iff $I \subseteq \sqrt{0}$.

Now suppose $\text{Spec}(A)$ is irreducible, meaning that if $\text{Spec}(A) =
V(I) \cup V(J) = V(IJ)$, then either $\text{Spec}(A) = V(I)$ or
$\text{Spec}(A)=V(J)$. This is equivalent to saying that whenever $IJ \subseteq
\sqrt{0}$ then $I \subseteq \sqrt{0}$ or $J \subseteq \sqrt{0}$, but this
exactly means $\sqrt{0}$ is a prime ideal.

Now assume $\sqrt{0} \subseteq A$ is a prime ideal and assume $\text{Spec}(A) =
V(I) \cup V(J) = V(IJ)$. Thus, $IJ \subseteq \sqrt{0}$, but since $\sqrt{0}$ is
prime, this implies either $I \subseteq \sqrt{0}$ or $J \subseteq \sqrt{0}$,
which is equivalent to saying $V(I)=\text{Spec}(A)$ or $V(J)=\text{Spec}(A)$.
Thus, $A$ is irreducible. \qedwhite

\begin{problem}[Exercise 4]
    Let $X$ be a topological space. Show the following:
    \begin{enumerate}
        \item[\textit{(a)}] If $Y \subseteq X$ is an irreducible subspace, then
        its closure $\overline{Y}$ is irreducible.
        \item[\textit{(b)}] Every irreducible subspace of $X$ is contained 
        in a maximal irreducible subspace.
        \item[\textit{(c)}] The maximal irreducible subspaces of $X$ are closed
        and cover $X$. They are called the \textbf{irreducible components} of
        $X$. What are the irreducible components of a Hausdorff space?
        \item[\textit{(d)}] If $A$ is a ring and $X=\text{Spec}(A)$, then the
        irreducible components of $X$ are the closed sets $V(\mathfrak{p})$,
        where $\mathfrak{p}$ is a minimal prime ideal of $A$.
    \end{enumerate}
\end{problem}

\textbf{Solution.}

\textit{Part (a).} Suppose that $\overline{Y}$ is reducible. Thus, there exist
closed subsets $G_0, G_1 \subsetneq \overline{Y}$ with $\overline{Y} = G_0 \cup
G_1$. Since $G_0$ and $G_1$ are closed in $\overline{Y}$, there exist closed
sets $X_0, X_1 \subseteq X$ such that $G_0 = X_0 \cap \overline{Y}$ and $G_1 =
X_1 \cap \overline{Y}$. Then, 
\begin{equation*}
    Y \subseteq \overline{Y} = G_0 \cup G_1 = (X_0 \cap \overline{Y}) \cup (X_1 \cap \overline{Y})
\end{equation*}

Therefore, $Y = (Y \cap X_0) \cup (Y \cap X_1)$, where $Y \cap X_i$ are closed
in $Y$. Since $Y$ is irreducible, one of these equals $Y$. Suppose $Y \cap X_0 =
Y$ or $Y \subseteq X_0$. But then $\overline{Y} \subseteq \overline{X}_0 = X_0$
so $\overline{Y} = \overline{Y} \cap A = G_0$. This is a contradiction. 

\textit{Part (b).} Suppose $\{A_{\alpha}\}_{\alpha \in I}$ be a chain of
irreducible subsets of $X$. We shall show that their union $\mathsf{A} :=
\bigcup_{\alpha \in I}A_{\alpha}$ is irreducible. Suppose it is not. Then, there
exist two closed sets $U,V \subsetneq \mathsf{A}$ such that $\mathsf{A} = U \cup
V$. Clearly, there exist two indices, say $\beta,\gamma$, such that $U \cap
A_{\beta}$ and $V \cap A_{\gamma}$ are non-empty. Since $\{A_{\alpha}\}_{\alpha
\in I}$ is a chain, either $A_{\beta} \subseteq A_{\gamma}$ or $A_{\gamma}
\subseteq A_{\beta}$. Suppose the former. Then, $U \cap A_{\gamma}$ and $V \cap
A_{\gamma}$ are both closed in $A_{\gamma}$, so either $U \cap
A_{\gamma}=A_{\gamma}$ or $V \cap A_{\gamma}=A_{\gamma}$ since $A_{\gamma}$ is
irreducible. This implies that either $A_{\gamma} \subseteq U$ or $A_{\gamma}
\subseteq V$. However, this is a contradiction: if $A_{\gamma} \subseteq U$,
then $V \cap A_{\gamma} = \emptyset$, hence the assumption of $\mathsf{A}$ being
irreducible is impossible. 

Now we simply apply the \textit{Zorn's lemma}: for each $Y \subseteq X$,
consider the family of irreducible sets that contain $Y$. The union of such
sets, as shown above, is clearly its maximal element. 

\textit{Part (c).} From \textit{part (b)}, every irreducible $Y$ is contained in
the corresponding maximal irreducible subspace $\mathsf{Y}$. By \textit{part
(a)}, we know that $\overline{\mathsf{Y}}$ is also irreducible. But since
$\mathsf{Y} \subseteq \overline{\mathsf{Y}}$ and $\mathsf{Y}$ is maximal, this
forces $\mathsf{Y} = \overline{\mathsf{Y}}$, which means $\mathsf{Y}$ is closed. 

Now we show that these maximal subspaces cover $X$. Take some $\{x\} \subseteq
X$. Clearly, $\{x\}$ is irreducible. From \textit{part (b)}, there exists some
maximal irreducible space $V_x$ such that $\{x\} \subseteq V_x$. Therefore,
$\bigcup_{x \in X}V_x$ covers $X$.

For Hausdorff space, irreducible components are simply singleton sets $\{x\}$.

\textit{Part (d).} Following Exercise 3, we know (with minor adjustments) that
$V(I)$ is irreducible iff $\sqrt{I}$ is a prime ideal. Since $V(I)=V(\sqrt{I})$,
this implies that a closed subset of $\text{Spec}(A)$ is closed iff it is of
form $V(\mathfrak{p})$ for some prime $\mathfrak{p} \subseteq A$.

At this point, we are almost done. First note that $V(\mathfrak{p})$ is an
irreducible component if $\mathfrak{p}$ is minimal. Firstly, it is clearly
irreducible. Suppose $U$ is an irreducible closed set such that $V(\mathfrak{p})
\subseteq U = V(\mathfrak{q})$. Then, $\mathfrak{q} \subseteq \mathfrak{p}$, but
since $\mathfrak{p}$ is minimal, $\mathfrak{q}=\mathfrak{p}$. Thus
$V(\mathfrak{p})$ is maximal. 

Finally, we have to show that every irreducible component is of form
$V(\mathfrak{p})$ for minimal $\mathfrak{p}$. But since every irreducible
component is closed, it is of form $V(\mathfrak{p})$ for some prime
$\mathfrak{p}$. For it to be maximal, as shown above, $\mathfrak{p}$ must be
minimal. \qedwhite

\begin{problem}[Exercise 5]
    Let $X$ be a topological space. We say that $X$ is
    \begin{enumerate}
        \item \textbf{sober} if every irreducible closed subset $Y \subseteq X$
        has a unique generic point, i.e. there exists a unique $\eta_Y \in Y$
        such that $Y=\overline{\{\eta_Y\}}$;
        \item \textbf{quasi-compact} (qc) if every open cover has a finite
        subcover;
        \item \textbf{spectral} if it is sober, qc, if the intersection of every
        two qc open subsets of X is qc, and if its qc open subsets form a base
        for the topology.
    \end{enumerate}
    Let $A$ be a commutative ring. Prove that $\text{Spec}(A)$ is a spectral
    space.
\end{problem}

\textbf{Proof.} $\text{Spec}(A)$ is clearly sober: every irreducible closed
subset $Y$ is of form $Y = V(\mathfrak{p})$, as we have shown in Exercise 3 (and
used in Exercise 4). But in Exercise 2 we have also shown that $V(\mathfrak{p})
= \overline{\{\eta_Y\}}$ for some $\eta_Y \in Y$, so $Y =
\overline{\{\eta_Y\}}$.

Now recall that open sets in $\text{Spec}(A)$ are $D(f) = \{\mathfrak{p} \in
\text{Spec}(A): f \notin \mathfrak{p}\}$. Suppose $\text{Spec}(A) =
\bigcup_{\alpha \in I} D(f_{\alpha})$ is an open cover.

Now, consider the ideal $\mathfrak{a} := \langle f_{\alpha} \rangle_{\alpha \in
I} \subseteq A$. Assume $\mathfrak{a} \subsetneq A$, then $\mathfrak{a}$ is
contained in some maximal ideal, and thus in some prime ideal $\mathfrak{p}$.
Note that for each $f_{\alpha}$, $f_{\alpha} \in \mathfrak{a} \subseteq
\mathfrak{p}$, so $\mathfrak{p} \not\in D(f_{\alpha})$, and thus $\mathfrak{p}
\not\in \bigcup_{\alpha \in I}D(f_{\alpha})$. But $\bigcup_{\alpha \in
I}D(f_{\alpha})$ covers the whole $\text{Spec}(A)$, so we reached the
contradiction. Thus, $\langle f_{\alpha} \rangle_{\alpha \in I} = A$, so in
particular there is a finite subset $J \subseteq I$ such that $\sum_{\alpha \in
J} c_{\alpha}f_{\alpha} = 1$. It is thus easy to see that $\bigcup_{\alpha \in
J}D(f_{\alpha}) = \text{Spec}(A)$ is a finite open subcover. 

Two left properties are then trivial: note that $D(f)$ is open, qc and forms a
basis for topology. And since $D(f) \cap D(g) = D(fg)$, it follows that
intersection of two qc open sets is again a qc open set. \qedwhite

\end{document}